%%%
%
% $Autor: Wings $
% $Datum: 2021-05-14 $
% $Pfad: ArduinoLab $
% $Dateiname: 
% $Version: 4620 $
%
% !TeX spellcheck = de_DE/GB
% !TeX program = pdflatex
% !BIB program = biber/bibtex
% !TeX encoding = utf8
%
%%%


\chapter{Arduino Lab for MicroPython}

\section{Einführung}

Arduino Lab for MicroPython ist ein Editor für MicroPython-Programme. Arduino Lab for MicroPython unterstützt die Verbindung mit einem Board, das Hochladen von Code, den Dateitransfer und eine interaktive REPL-Shell.

\section{Installation}
Die Schritte zur Konfiguration von Arduino Alvik für Arduino Lab for MicroPython sind wie folgt:
\begin{enumerate}
	\item Stellen den Ein/Aus-Schalter des Arduino Alvik Roboters auf ``Aus''. 
	\item Verbinden den ESP32 auf dem Alvik über ein USB-C-Kabel mit dem Computer.
	\item Laden Arduino Lab for MicroPython von der offiziellen Arduino Website herunter.
	\item Greift über \menu{File Manager} auf den Ordnerinhalt zu oder wechselt zu einem neuen Ordner. Wie in der Abbildung \ref{Filemanager} dargestellt \cite{Jose:2025}.	
	
	\begin{figure}[H]  
		\begin{center}
			\includegraphics[width=12cm]{Arduino/Alvik/Filemanager}
			\caption{Arduino lab for MicroPython File manager} 
			\label{Filemanager}                                                                                                                                                                                                                                                                                                          
		\end{center}
	\end{figure}
	\item Klicken auf die Schaltfläche \menu{Connect} und wählen den richtigen COM-Anschluss. Wie in der Abbildung \ref{Arduinolabconnect} dargestellt \cite{Jose:2025}.
	\begin{figure}[H]  
		\begin{center}
			\includegraphics[width=12cm]{Arduino/Alvik/Arduinolabconnect}
			\caption{Arduino lab for MicroPython Connect} 
			\label{Arduinolabconnect}                                                                                                                                                                                                                                                                                                          
		\end{center}
	\end{figure}
	\newpage
	\item Klicken auf die Schaltflächen \menu{Editor und REPL}. Der obere wei\ss e Teil ist der Editor zum Schreiben von Skripten und der untere schwarze Teil ist das Terminal. Wie in der Abbildung \ref{EditorandREPL} dargestellt \cite{Jose:2025}.
	\begin{figure}[H]  
		\begin{center}
			\includegraphics[width=12cm]{Arduino/Alvik/EditorAndREPL}
			\caption{Editor und REPL} 
			\label{EditorandREPL}                                                                                                                                                                                                                                                                                                          
		\end{center}
	\end{figure}
	
	\Mynote{Schwarzbereich aus Bild entfernen und überall newpage entfernen}
	\newpage
	\item Klicken auf die Schaltfläche \menu{Run}, um das Standard-Beispielskript auszuführen. Nachdem das Skript ausgeführt wurde, wird ``Hallo, MicroPython!'' auf dem Terminal ausgegeben. Wie in der Abbildung \ref{Run} dargestellt \cite{Jose:2025}. 
	\begin{figure}[H]  
		\begin{center}
			\includegraphics[width=12cm]{Arduino/Alvik/Run}
			\caption{Arduino lab for MicroPython Run} 
			\label{Run}                                                                                                                                                                                                                                                                                                          
		\end{center}
	\end{figure}
	\item Öffnen den \menu{File Manager} eine neue Datei. Benennen und öffnen die neue Datei.
	\item Schieben den Ein/Aus-Schalter des Alvik in die Position ``Ein'' \cite{Jose:2025}. 
	
\end{enumerate}

\newpage
\section{Erste Schritte}
Alvik kann zunächst programmiert werden, indem die Text ein- und ausgegeben werden. Der Code \ref{Eingabetext} ist wie folgt:
%\begin{figure}[ht]
%	\begin{center}
%		\includegraphics[width=12cm]{Arduino/Alvik/Input}
%		\caption{Eingabe und Ausgabe von Text Code} 
%		\label{Eingabetext}                                                                                                                                                                                                                                                                                                          
%	\end{center}
%\end{figure}

	{
	\captionof{code}{Eingabe und Ausgabe von Text Code}\label{Eingabetext}	
	{	\pythonstyle
		\lstinputlisting[]{../../Code/Alvik/Communikation/Textausgeben.py}}
}
\begin{enumerate}
	\item \textbf{Eingabe und Ausgabe von Text}
	\begin{itemize}
		\item \textbf{Importieren Module:} Importieren von sleep aus dem Zeitmodul und Verwendung der Funktion \PYTHON{sleep()}, um das Programm einige Sekunden warten zu lassen, bevor es mit der nächsten Codezeile fortfährt.
		\item \textbf{Eingabe- und Ausgabezeichen:} Die Funktion \PYTHON{input()} wird zur Eingabe von Informationen in das Terminal verwendet. Alle vom Benutzer eingegebenen Informationen werden in der Variablen student\_name gespeichert. Hier wartet das Programm auf die Eingabe und fährt erst dann fort, wenn die Eingabe eingegangen ist. Verwenden \PYTHON{print()}, um Text auf dem Terminal auszugeben.
		\item \textbf{Verkettung:} Verwenden den \PYTHON{+} Symbol, um die Zeichenkette mit der Variable \PYTHON{student\_name} zu kombinieren.
		\item \textbf{Einstellung der Interaktionsgeschwindigkeit:} Verwenden von \PYTHON{sleep()}, um die Interaktion zu verlangsamen, was zu einem natürlicheren Dialogfluss führt.
		
	\end{itemize}
\end{enumerate}	

\section{Datei im Arduino Alvik speichern}

Arduino Alvik ermöglicht das Ausführen integrierter Programme, indem Dateien vom Computer über Arduino Lab für MicroPython in den Arduino Alvik importiert werden.
	\newpage
Die Schritte für die Übertragung von Dokumenten sind wie folgt:
\begin{itemize}

	\item Öffnen die Arduino Lab for MicroPython-Datei über die Schaltfläche \menu{File Manager}. Wie in der Abbildung \ref{alab} gezeigt, werden im linken Fenster die in Arduino Alvik integrierten Dateien und im rechten Fenster die auf dem Computer gespeicherten Dateien angezeigt.
	\begin{figure}[H]
		\begin{center}
			\includegraphics[width=12cm]{Arduino/Alvik/alab}
			\caption{Arduino Lab for MicroPython Dateispeicherort} 
			\label{alab}                                                                                                                                                                                                                                                                                                          
		\end{center}
	\end{figure}
	\newpage
	\item Wählen ein Programm aus und klicken auf die Schaltfläche \menu{Hochladen} (siehe Abb. \ref{hochladen}) , um das Programm von dem Computer auf das Arduino Alvik hochzuladen.
	\begin{figure}[H]
		\begin{center}
			\includegraphics[width=12cm]{Arduino/Alvik/Upload}
			\caption{Schaltfläche ``Hochladen''} 
			\label{hochladen}                                                                                                                                                                                                                                                                                                          
		\end{center}
	\end{figure}
	\item Öffnen die in Arduino Alvik integrierte Datei \FILE{main.py} \ref{main}, löschen den darin enthaltenen Code, geben den Code \PYTHON{import xxx} ein (wobei xxx der Name der zu importierenden Datei ist) und speichern anschlie\ss end.
	
	\begin{Python} [caption={Code für main.py},label={main}]
		import xxx #(wobei xxx der Name der zu importierenden Datei ist)
	\end{Python}
	
	\item Trennen das Arduino Alvik vom Computer und schalten es ein. Das Arduino Alvik führt automatisch die zuvor eingerichteten eingebauten Programme aus.
	\item Beenden das Programm durch Ausschalten des Schalters am Arduino Alvik \cite{Paolo:2025}. 
\end{itemize}



